\documentclass[11pt]{article}
\usepackage[a4paper,margin=27mm]{geometry}
\usepackage{amsmath,amssymb,amsthm,booktabs,array}
\usepackage[T1]{fontenc}
\usepackage{lmodern}
\usepackage[hidelinks,breaklinks]{hyperref}
\setlength{\emergencystretch}{2em}
\newtheorem{theorem}{Theorem}
\newtheorem{lemma}[theorem]{Lemma}
\newtheorem{corollary}[theorem]{Corollary}
\newcommand{\Z}{\mathbb Z}
\newcommand{\Q}{\mathbb Q}
\newcommand{\vp}{v_p}
\title{Proof of Hanna's conjecture in OEIS A181080,\\COMMENTS, first line:\\integrality and binomial-power companions}
\author{Haobo Ma\\\small ChronoAI Pte Ltd, Singapore
\and Wenlin Zhang\\\small National University of Singapore\\\small ChronoAI Pte Ltd, Singapore}
\date{8 November 2026}
\begin{document}
\maketitle
\begin{abstract}
We prove Paul D. Hanna's integrality conjecture in OEIS A181080,
COMMENTS, first line (2 October 2010), for the series specified in NAME:
\[
 \exp\!\left(\sum_{n\geq1}\sum_{k=0}^n\binom nk^{n-k+1}\frac{x^{n+k}}n\right)
\]
has integer coefficients.
Prime-scaling congruences for binomial coefficients imply Dold congruences for
its logarithmic derivative. The same argument proves the related assertions
for A181070, A181074, A181076, A181078, A181082, A181084 and A166894,
and gives integer logarithmic sequences in all eight cases. A linear-exponent
extension applies to powers of any nonnegative integral series with constant
term one. We also verify Hanna's three nonintegral comparison series and
correct the exponential cross-reference in A181079.
\end{abstract}

\section{The conjecture and the integrality criterion}
Hanna's entry A181080~\cite{H80} states, ``Conjecture: this sequence consists
entirely of integers.'' Its NAME specifies the formal series
\begin{equation}\label{eq:main}
 A(x)=\exp L(x),\qquad
 L(x)=\sum_{n\geq1}\sum_{k=0}^n\frac{\binom nk^{n-k+1}}{n}x^{n+k}
 \quad\text{in }x\Q[[x]].
\end{equation}
Every coefficient is a finite sum, so all operations below are formal;
analytic convergence is unnecessary.
\begin{theorem}\label{thm:main}
The series $A(x)$ in \eqref{eq:main} belongs to $1+x\Z[[x]]$.
\end{theorem}
The relevant standard framework is the equivalence between Dold congruences,
integral M\"obius transforms and integral Euler products; see
Byszewski, Graff and Ward~\cite[Definition 2.1, Theorem 2.5 and Lemma 2.6]{BGW}.
Moree's Witt-transform theorem~\cite[Theorems 2 and 4(2)]{Moree}
is an earlier formal-series integrality result. These general frameworks
provide the Euler-product mechanism; the variable-exponent congruences
needed for Hanna's arrays are established below.
We include the needed direction
to specify exactly how the congruences give integrality.
For a prime $p$ and positive integer $n$, $\vp(n)$ denotes the exponent of $p$ in $n$.

\begin{lemma}\label{lem:dold}
Suppose integers $c_m$, $m\geq1$, satisfy
\begin{equation}\label{eq:dold}
 c_{pm}\equiv c_m\pmod{p^{\vp(m)+1}}
 \qquad(p\text{ prime},\ m\geq1).
\end{equation}
Then $\exp(\sum_{m\geq1}c_mx^m/m)$ has integer coefficients.
\end{lemma}
\begin{proof}
Put $b_j=j^{-1}\sum_{d\mid j}\mu(d)c_{j/d}$, where $\mu$ is the M\"obius function.
For $j=p^a u$, $a\geq1$, $p\nmid u$, the numerator is
\[
 \sum_{d\mid u}\mu(d)\bigl(c_{p^a u/d}-c_{p^{a-1}u/d}\bigr),
\]
which is divisible by $p^a$ by \eqref{eq:dold}. This holds for every prime
power dividing $j$, hence $b_j\in\Z$.
M\"obius inversion gives $c_m=\sum_{j\mid m}j b_j$, and therefore
\[
 \exp\!\left(\sum_{m\geq1}\frac{c_m}{m}x^m\right)
   =\prod_{j\geq1}(1-x^j)^{-b_j}\in1+x\Z[[x]].
\]
Each factor has integral coefficients for every integer $b_j$; the product
is well defined coefficient by coefficient.
\end{proof}

\section{Two prime-scaling lemmas}
\begin{lemma}\label{lem:coeff}
For $U(t)\in\Z[[t]]$, write $B(n,k)=[t^k]U(t)^n$. For $n\geq1$, $k\geq0$
and every prime $p$,
\begin{equation}\label{eq:coeff}
 n\mid kB(n,k),\qquad
 B(pn,pk)\equiv B(n,k)\pmod{p^{\vp(n)+1}}.
\end{equation}
In particular, these statements hold for $B(n,k)=\binom nk$ and
$B(n,k)=\binom{n+k-1}{k}$.
\end{lemma}
\begin{proof}
The identity $(U^n)'=nU^{n-1}U'$ gives the divisibility by taking the
coefficient of $t^{k-1}$; for $k=0$ it is immediate.
The Frobenius identity modulo $p$ gives
$U(t)^p\equiv U(t^p)\pmod p$.
If $R\equiv S\pmod{p^s}$, $s\geq1$, the binomial theorem implies
$R^p\equiv S^p\pmod{p^{s+1}}$: the intermediate terms acquire an extra
factor $p$, and $sp\geq s+1$ for the last term.
Write $n=p^a u$, $p\nmid u$. Applying this observation $a$ times and then
raising to the $u$th power yields
$U(t)^{pn}\equiv U(t^p)^n\pmod{p^{a+1}}$.
Taking coefficients of $t^{pk}$ proves \eqref{eq:coeff}.
The final assertions use $U=1+t$ and $U=(1-t)^{-1}$, respectively.
\end{proof}

\begin{lemma}\label{lem:exponent}
Let $n\geq1$, $b,t\geq0$ be integers with $n\mid tb$. Then
\begin{equation}\label{eq:exponent}
 b^{pt+1}\equiv b^{t+1}\pmod{p^{\vp(n)+1}}
 \quad(p\text{ prime}).
\end{equation}
Also, if $e\geq n$ and $n\mid eb$, then
\begin{equation}\label{eq:largeexponent}
 b^{pe}\equiv b^e\pmod{p^{\vp(n)+1}}.
\end{equation}
\end{lemma}
\begin{proof}
Write $a=\vp(n)$. For \eqref{eq:exponent}, the cases $t=0$ and $b=0$ are
immediate. If $p\nmid b$, the divisibility gives $p^a\mid t$.
Euler's theorem modulo $p^{a+1}$ applies, since the two exponents differ
by $(p-1)t$, a multiple of $p^a(p-1)$.
If $p\mid b$, put $s=\vp(b)\geq1$. Then
$a\leq\vp(t)+s$, while $t\geq\vp(t)+1$ implies
$(t+1)s\geq\vp(t)+s+1\geq a+1$.
Both powers in \eqref{eq:exponent} are thus divisible by $p^{a+1}$.
For \eqref{eq:largeexponent}, $b=0$ is immediate. In the unit case,
$p^a\mid e$ and the same Euler argument applies.
Otherwise $e\vp(b)\geq n\geq p^a\geq a+1$, so both powers again vanish
modulo $p^{a+1}$. The inequality $p^a\geq a+1$ follows by induction on $a$.
\end{proof}

\section{From an array to an integral exponential}
\begin{lemma}\label{lem:array}
Let $F(n,k)\in\Z$ for $n\geq1$, $k\geq0$. Suppose
\begin{equation}\label{eq:array}
 n\mid kF(n,k),\qquad
 F(pn,pk)\equiv F(n,k)\pmod{p^{\vp(n)+1}}
\end{equation}
for every prime $p$. Then
$\exp(\sum_{n\geq1,k\geq0}F(n,k)x^{n+k}/n)\in1+x\Z[[x]]$.
\end{lemma}
\begin{proof}
The first condition makes
$T(n,k)=(n+k)F(n,k)/n$ an integer. Define
\[
 c_m=\sum_{n=1}^m T(n,m-n).
\]
These are exactly $m$ times the logarithm coefficients.
Fix $p$ and $m\geq1$. In $c_{pm}$, terms whose first index $n$ is not
divisible by $p$ vanish modulo $p^{\vp(m)+1}$: indeed
$nT(n,pm-n)=pmF(n,pm-n)$ and $n$ is a unit modulo that power of $p$.
Pair each remaining term with $T(n,m-n)$ in $c_m$. Their difference is
\[
 T(pn,p(m-n))-T(n,m-n)
   =\frac mn\bigl(F(pn,p(m-n))-F(n,m-n)\bigr).
\]
By the second condition its $p$-adic valuation is at least
$\vp(m)-\vp(n)+(\vp(n)+1)=\vp(m)+1$.
The displayed difference is an integer; thus this valuation bound is an
ordinary integer divisibility statement. A zero difference satisfies it as well.
Summing proves \eqref{eq:dold}, so Lemma~\ref{lem:dold} applies.
\end{proof}

\begin{proof}[Proof of Theorem~\ref{thm:main}]
For $0\leq k\leq n$, put $B=\binom nk$ and $F(n,k)=B^{n-k+1}$;
put $F(n,k)=0$ for $k>n$. Since $n\mid kB$, we have $n\mid kF(n,k)$.
Also $n\mid(n-k)B$, because $n\mid nB-kB$.
For a prime $p$, set $q=p^{\vp(n)+1}$. Lemma~\ref{lem:coeff} gives
$\binom{pn}{pk}\equiv B\pmod q$, so for $k\leq n$,
\[
 F(pn,pk)\equiv B^{p(n-k)+1}
   \equiv B^{n-k+1}=F(n,k)\pmod q,
\]
where the second step uses Lemma~\ref{lem:exponent} with $t=n-k$.
For $k>n$ both sides are zero. Lemma~\ref{lem:array} now proves the theorem.
\end{proof}

\section{Linear exponents and companion entries}
\begin{corollary}\label{cor:linear}
Let $U(t)\in\Z_{\geq0}[[t]]$, $U(0)=1$, and put
$B(n,k)=[t^k]U(t)^n$. The series
\[
 \exp\!\left(\sum_{n\geq1,k\geq0}
 B(n,k)^{rn+sk+\delta}\frac{x^{n+k}}n\right)
\]
has integer coefficients in each of the following ranges of integer parameters:
\[
\begin{array}{c|c}
\delta\geq1 & r,s\geq0\\
\delta=0 & r\geq1,\ s\geq0\\
\delta=-1 & r,s\geq1.
\end{array}
\]
All exponents are nonnegative. In the third range the exponent can vanish
only at $n=1,k=0,r=1$, where $B(n,0)=1$.
\end{corollary}
\begin{proof}
Write $u=rn+sk$ and $B=B(n,k)$. Lemma~\ref{lem:coeff} gives
$n\mid uB$. For $\delta\geq1$, multiply \eqref{eq:exponent}, with
$t=u$, by $B^{\delta-1}$. For $\delta=0$, use
\eqref{eq:largeexponent} with $e=u\geq n$.
In both cases, first replace $B(pn,pk)$ by $B$ modulo
$p^{\vp(n)+1}$ using Lemma~\ref{lem:coeff}.
For $\delta=-1$, $k=0$ gives $B=1$ and is immediate. If $k>0$, then
$u-1\geq n$. When $p\mid B$ and $B>0$, both
$B^{pu-1}$ and $B^{u-1}$ vanish modulo $p^{\vp(n)+1}$,
since $n\geq p^{\vp(n)}\geq\vp(n)+1$; when $B=0$ both are zero.
When $p\nmid B$, $n\mid uB$ implies $p^{\vp(n)}\mid u$,
so Euler's theorem applies to the exponent difference $(p-1)u$.
Thus the second condition of \eqref{eq:array} holds in all three ranges.
The first follows from $n\mid kB$, because the exponent is positive
unless $k=0$. Lemma~\ref{lem:array} proves the claim.
\end{proof}

\begin{corollary}\label{cor:siblings}
The seven companion exponentials in Table~\ref{tab:families} have
integer coefficients. In each of its eight rows the sequence
$c_m=m[x^m]\log A(x)$ is the indicated integer sequence and satisfies
\eqref{eq:dold}.
\end{corollary}
\begin{proof}
For A181070~\cite{H70}, take $B=\binom nk$ and $F=B^{k+1}$,
zero outside $k\leq n$. Use \eqref{eq:exponent} with $t=k$.
For A181074~\cite{H74}, use $B=\binom{n+k-1}{k}$ and $F=B^{k+1}$
for all $k\geq0$; Lemma~\ref{lem:coeff} applies with $U=(1-t)^{-1}$,
and the same exponent argument proves \eqref{eq:array}.
For A166894~\cite{H894} and A181082~\cite{H82}, use $B=\binom nk$
with exponents $e=n$ and $e=n+k$, respectively, and set $F=0$ when $k>n$.
Both exponents satisfy $e\geq n$ and $n\mid eB$, so
\eqref{eq:largeexponent} and Lemma~\ref{lem:coeff} give the second
condition of \eqref{eq:array}. The first follows from $n\mid kB$ and $e\geq1$.
Lemma~\ref{lem:array} proves these four assertions.
For A181076 and A181078~\cite{H76,H78}, use
$B=\binom{n+k-1}{k}$ with exponents $n$ and $n+k-1$, respectively;
Corollary~\ref{cor:linear} applies with $(r,s,\delta)=(1,0,0)$
and $(1,1,-1)$. For A181084~\cite{H84}, use $B=\binom nk$
and $(r,s,\delta)=(1,1,1)$, with zero values for $k>n$.
Collecting terms of degree $m$ in \eqref{eq:main} gives
\begin{equation}\label{eq:ghost}
 c_m=\sum_{k=0}^{\lfloor m/2\rfloor}
       \frac{m}{m-k}\binom{m-k}{k}^{m-2k+1},
\end{equation}
the logarithm formula of A181081~\cite{H81}. For every other row the same
substitution $n=m-k$ gives
\[
 c_m=\sum_k\frac{m}{m-k}\,B(m-k,k)^{e(m-k,k)}.
\]
For binomial rows the range is $0\leq k\leq\lfloor m/2\rfloor$;
for negative-binomial rows it is $0\leq k\leq m-1$, and their base becomes
$\binom{m-1}{k}$. These are exactly the NAME or L.g.f. formulas of the
logarithmic entries in Table~\ref{tab:families}~\cite{Hlogs}.
Lemma~\ref{lem:array} proves their integrality and Dold congruences.
\end{proof}

\begin{table}[ht]
\centering\small
\renewcommand{\arraystretch}{1.65}
\begin{tabular}{@{}llll@{}}\toprule
Exponential & $B(n,k)$ & $e(n,k)$ & $c_m$ entry\\\midrule
A181080 & $\binom nk$ & $n-k+1$ & A181081\\
A181070 & $\binom nk$ & $k+1$ & A181071\\
A166894 & $\binom nk$ & $n$ & A166895\\
A181082 & $\binom nk$ & $n+k$ & A181083\\
A181084 & $\binom nk$ & $n+k+1$ & A181085\\
A181074 & $\binom{n+k-1}{k}$ & $k+1$ & A181075\\
A181076 & $\binom{n+k-1}{k}$ & $n$ & A181077\\
A181078 & $\binom{n+k-1}{k}$ & $n+k-1$ & A181079\\\bottomrule
\end{tabular}
\caption{In every row, $A(x)=\exp(\sum B(n,k)^{e(n,k)}x^{n+k}/n)$
and $c_m=m[x^m]\log A(x)$. Binomial rows have $0\leq k\leq n$;
negative-binomial rows have $k\geq0$.}\label{tab:families}
\end{table}

\section{Comparison series and source coverage}
The COMMENTS fields of A181080, A181070 and A181074 distinguish series with
the exponent decreased by one. We confirm their stated nonintegrality,
rather than contradicting an integrality assertion about them.
For $\widetilde A=\exp(\sum_{m\geq1}\widetilde c_m x^m/m)
=\sum_{m\geq0}\widetilde a_m x^m$, coefficient comparison in
$x\widetilde A'=\widetilde A\sum_{m\geq1}\widetilde c_m x^m$ gives
\[
 \widetilde a_0=1,\qquad
 m\widetilde a_m=\sum_{j=1}^m\widetilde c_j\widetilde a_{m-j}.
\]
Direct finite sums from the three specified logarithms give the following
exact values; the recurrence then gives the last column.
\begin{center}\small
\begin{tabular}{@{}lll@{}}\toprule
Comparison in & $(\widetilde c_1,\ldots,\widetilde c_6)$ & $\widetilde a_6$\\\midrule
A181080, exponent $n-k$ & $(1,3,4,15,86,807)$ & $327/2$\\
A181070, exponent $k$ & $(1,3,4,7,21,63)$ & $45/2$\\
A181074, exponent $k$ & $(1,3,7,27,231,4038)$ & $1481/2$\\\bottomrule
\end{tabular}
\end{center}
For each row, degrees zero through five are integral, so degree six is the
first obstruction. The convention $0^0$ does not arise in the bounded sums,
since $\binom nk>0$ for $0\leq k\leq n$.

There is one cross-reference to correct. The FORMULA field of A181079
states, ``Logarithmic derivative of A181076.'' The correct exponential is
A181078, as its L.g.f. and EXAMPLE already specify. Indeed, Table~\ref{tab:families}
gives $xA'_{181078}(x)/A_{181078}(x)=\sum_{m\geq1}c_mx^m$
with $c_4=95$, whereas $xA'_{181076}(x)/A_{181076}(x)$ has coefficient
$59$ at $x^4$ (A181077). This both disproves the printed identification
and supplies its corrected identity. In A181081, the typographical
\texttt{n-2k*+1} in NAME is read as $n-2k+1$ by its L.g.f. and EXAMPLE,
which agree with \eqref{eq:ghost}.

Hanna's first COMMENTS line in each of A181070, A181074, A181076,
A181078, A181082 and A181084 repeats the integrality conjecture quoted
in Section~1. A166894 instead asserts ``an integer series in x'' in NAME.
The following table separates these source assertions from additional
consequences and from unrelated or uncovered items.
\begin{center}\small
\begin{tabular}{@{}>{\raggedright\arraybackslash}p{.31\linewidth}>{\raggedright\arraybackslash}p{.64\linewidth}@{}}\toprule
Source and exact locator & Coverage\\\midrule
A181080, NAME and first COMMENTS & Proved: integrality, Theorem~\ref{thm:main}.\\
A181070, A181074, A181076, A181078, A181082, A181084; NAME and first COMMENTS
 & Proved: the six integrality conjectures, Corollary~\ref{cor:siblings}.\\
A166894, NAME & Proved: the asserted integrality.\\
Eight logarithmic entries in Table~\ref{tab:families}, NAME and L.g.f./EXAMPLE
 & Formulas verified; integrality proved. Dold congruences are not posed;
 proved here for all eight.\\
A181079, FORMULA, second sentence & Refuted-and-corrected: replace A181076 by
A181078; coefficients $59\ne95$ at degree four.\\
A181080, A181070, A181074; COMMENTS, comparison formulas
 & Proved: the stated nonintegrality, witnessed at degree six.\\
Corollary~\ref{cor:linear} & Not posed; proved: the linear-exponent families.\\
A166895, FORMULA, limit (Kotesovec, 20 November 2024)
 & Not covered: the separate asymptotic formula; it is not an integrality question.\\
A181072 and A181073, NAME & Not covered: unrelated random-walk and permutation
entries appearing among variant cross-references~\cite{Hmisc}; no identification
with these exponentials is asserted.\\\bottomrule
\end{tabular}
\end{center}

\section{Mathematical and formal scope}
The proofs above establish all covered statements without finite premises.
The pinned Lean declarations at revision \texttt{abb48f7} have the following
scopes; the common namespace prefix is
\texttt{D5.S1.Recurrence.Algebraic.}
\begin{itemize}
\raggedright
\item \href{https://github.com/the-omega-institute/trureturing/blob/abb48f723fe221da235c625800ba0d2ca1755296/D5/S1/Recurrence/Algebraic/HannaA181080Integrality.lean}{\texttt{HannaA181080Integrality.lean}}:
\texttt{HannaA181080Integrality.result} proves integrality in every degree
for the literal bounded rational logarithm in \eqref{eq:main} and its formal exponential.
\item \href{https://github.com/the-omega-institute/trureturing/blob/abb48f723fe221da235c625800ba0d2ca1755296/D5/S1/Recurrence/Algebraic/HannaA181070Integrality.lean}{\texttt{HannaA181070Integrality.lean}}:
\texttt{HannaA181070Integrality.result} proves the A181070 assertion;
\texttt{integralGeneric} proves the nonnegative-array version of
Lemma~\ref{lem:array} for the diagonal logarithm, and
\texttt{literalBounded}/\texttt{literalUnbounded} identify its coefficient formulas.
\item \href{https://github.com/the-omega-institute/trureturing/blob/abb48f723fe221da235c625800ba0d2ca1755296/D5/S1/Recurrence/Algebraic/HannaA181074Integrality.lean}{\texttt{HannaA181074Integrality.lean}}:
\texttt{HannaA181074Integrality.result} proves integrality for the literal
negative-binomial logarithm with exponent $k+1$.
\end{itemize}
These three named \texttt{result} declarations correspond to A181080,
A181070 and A181074, respectively. The other five exponentials, the
linear-exponent extension, all eight logarithmic consequences, and the
source corrections are covered by the written arguments above.

The accompanying \texttt{verify.py} checks coefficients through degree 48
for all eight integral exponentials by a differential recurrence and Euler
products, and through degree 16 by the direct exponential sum. It checks
the two array conditions for $1\leq n\leq24$, $0\leq k\leq24$ and
$p\in\{2,3,5,7,11,13\}$, and Dold congruences when $pm\leq48$ for these
primes. It also checks the eight logarithmic formulas and the linear-exponent
families on three integral series and seven parameter choices through degree 12,
with prime-scaling checks for $1\leq n\leq6$, $0\leq k\leq12$ and
$p\in\{2,3,5\}$. These are finite checks, distinct from the all-degree proofs.

\begin{thebibliography}{99}\small\raggedright
\bibitem{H80} P. D. Hanna, OEIS A181080, 2 October 2010,
NAME, COMMENTS and CROSSREFS, \url{https://oeis.org/A181080}.
\bibitem{H70} P. D. Hanna, OEIS A181070, 2 October 2010,
NAME and COMMENTS, \url{https://oeis.org/A181070}.
\bibitem{H74} P. D. Hanna, OEIS A181074, 2 October 2010,
NAME and COMMENTS, \url{https://oeis.org/A181074}.
\bibitem{H894} P. D. Hanna, OEIS A166894, 23 November 2009,
NAME, \url{https://oeis.org/A166894}.
\bibitem{H82} P. D. Hanna, OEIS A181082, 2 October 2010,
NAME and COMMENTS, \url{https://oeis.org/A181082}.
\bibitem{H81} P. D. Hanna, OEIS A181081, 2 October 2010,
FORMULA and examples, \url{https://oeis.org/A181081}.
\bibitem{H76} P. D. Hanna, OEIS A181076, 2 October 2010,
NAME, first COMMENTS and CROSSREFS, \url{https://oeis.org/A181076}.
\bibitem{H78} P. D. Hanna, OEIS A181078, 3 October 2010,
NAME, first COMMENTS and CROSSREFS, \url{https://oeis.org/A181078}.
\bibitem{H84} P. D. Hanna, OEIS A181084, dated 28 October 2010,
NAME and first COMMENTS, \url{https://oeis.org/A181084}.
\bibitem{Hlogs} P. D. Hanna, logarithmic companion entries, NAME and
L.g.f./EXAMPLE: \url{https://oeis.org/A166895},
\url{https://oeis.org/A181071}, \url{https://oeis.org/A181075},
\url{https://oeis.org/A181077}, \url{https://oeis.org/A181079},
\url{https://oeis.org/A181083}, \url{https://oeis.org/A181085}.
The separate limit formula in A166895 is attributed there to
V. Kotesovec, 20 November 2024.
\bibitem{Hmisc} OEIS A181072 and A181073, NAME,
\url{https://oeis.org/A181072}, \url{https://oeis.org/A181073}.
\bibitem{Moree} P. Moree, \emph{The formal series Witt transform},
Discrete Math. \textbf{295} (2005), 143--160,
\href{https://doi.org/10.1016/j.disc.2005.03.004}{doi:10.1016/j.disc.2005.03.004};
\href{https://arxiv.org/abs/math/0311194}{arXiv:math/0311194}, v2 (31 March 2005).
\bibitem{BGW} J. Byszewski, G. Graff and T. Ward,
\emph{Dold sequences, periodic points, and dynamics},
Bull. Lond. Math. Soc. \textbf{53} (2021), 1263--1298.
\href{https://doi.org/10.1112/blms.12531}{doi:10.1112/blms.12531};
\href{https://arxiv.org/abs/2007.04031}{arXiv:2007.04031}.
\end{thebibliography}
\bigskip\noindent\small
\copyright\ 2026 Haobo Ma and Wenlin Zhang.
Licensed under Creative Commons Attribution 4.0 International (CC BY 4.0).
\end{document}
